\section{Introduction}\label{sec:intro}
% 1. Clinical context: PAH, why PA haemodynamics (WSS, pulse wave, compliance) matter.
% 2. Pigs as the preclinical/xenotransplant model for the pulmonary circulation;
%    known anatomical differences (branching pattern, diameters, wall properties).
% 3. Computational PA haemodynamics so far: rigid-wall CFD vs FSI; mostly
%    single-species, small cohorts. Gap: no cohort-level FSI cross-species comparison.
% 4. Aim and contributions: automated, reproducible STL->FSI pipeline;
%    cohort comparison of human vs porcine; (healthy vs PAH?).
\todo{Introduction}
Large animal models are used widely in cardiovascular research and medical device development and certification.
Pigs and sheep are the most commonly used animals due to their physiological similarities to humans (body and heart weight, heart rate and ventricular kinetics) alongside their availability.
Additionally, the same factors have lead to pigs becoming prime candidates for cardiac xenotransplantation.

\todo{Brief description of cardiac xenotransplantation and potentially bearing of pulmonary arteries}

Hence, an understanding of the limits of the similarities between pulmonary haemodynamics in humans and pigs is important to assess the applicability of findings in these animal models to humans. 



These cases and example
